\bibitem{bobbin2024} M.~P. Bobbin, S. Sharlin, P. Feyzishendi, A.~H. Dang, C.~M. Wraback and T.~R. Josephson.
Formalizing chemical physics using the Lean theorem prover.
\emph{Digital Discovery} 3 (2024), 264--280.
\href{https://doi.org/10.1039/D3DD00077J}{doi:10.1039/D3DD00077J}.

\bibitem{butch2000} A.~W. Butch.
Dilution protocols for detection of hook effects/prozone phenomenon.
\emph{Clinical Chemistry} 46 (2000), 1719--1720.
\href{https://doi.org/10.1093/clinchem/46.10.1719}{doi:10.1093/clinchem/46.10.1719}.

\bibitem{demoura2021} L. de~Moura and S. Ullrich.
The Lean~4 theorem prover and programming language.
In \emph{Automated Deduction -- CADE 28}, Lecture Notes in Computer Science 12699, Springer, 2021, pp.~625--635.
\href{https://doi.org/10.1007/978-3-030-79876-5_37}{doi:10.1007/978-3-030-79876-5\_37}.

\bibitem{desilva2003} B. DeSilva, W. Smith, R. Weiner, M. Kelley, J. Smolec, B. Lee, M. Khan, R. Tacey, H. Hill and A. Celniker.
Recommendations for the bioanalytical method validation of ligand-binding assays to support pharmacokinetic assessments of macromolecules.
\emph{Pharmaceutical Research} 20 (2003), 1885--1900.
\href{https://doi.org/10.1023/B:PHAM.0000003390.51761.3d}{doi:10.1023/B:PHAM.0000003390.51761.3d}.

\bibitem{douglass2013} E.~F. Douglass~Jr., C.~J. Miller, G. Sparer, H. Shapiro and D.~A. Spiegel.
A comprehensive mathematical model for three-body binding equilibria.
\emph{Journal of the American Chemical Society} 135 (2013), 6092--6099.
\href{https://doi.org/10.1021/ja311795d}{doi:10.1021/ja311795d}.

\bibitem{dudley1985} R.~A. Dudley, P. Edwards, R.~P. Ekins, D.~J. Finney, I.~G. McKenzie, G.~M. Raab, D. Rodbard and R.~P. Rodgers.
Guidelines for immunoassay data processing.
\emph{Clinical Chemistry} 31 (1985), 1264--1271.
\href{https://doi.org/10.1093/clinchem/31.8.1264}{doi:10.1093/clinchem/31.8.1264}.

\bibitem{farouki2012} R.~T. Farouki.
The Bernstein polynomial basis: a centennial retrospective.
\emph{Computer Aided Geometric Design} 29 (2012), 379--419.
\href{https://doi.org/10.1016/j.cagd.2012.03.001}{doi:10.1016/j.cagd.2012.03.001}.

\bibitem{fernando1992a} S.~A. Fernando and G.~S. Wilson.
Studies of the `hook' effect in the one-step sandwich immunoassay.
\emph{Journal of Immunological Methods} 151 (1992), 47--66.
\href{https://doi.org/10.1016/0022-1759(92)90104-2}{doi:10.1016/0022-1759(92)90104-2}.

\bibitem{fernando1992b} S.~A. Fernando and G.~S. Wilson.
Multiple epitope interactions in the two-step sandwich immunoassay.
\emph{Journal of Immunological Methods} 151 (1992), 67--86.
\href{https://doi.org/10.1016/0022-1759(92)90105-3}{doi:10.1016/0022-1759(92)90105-3}.

\bibitem{garloff1993} J. Garloff.
The Bernstein algorithm.
\emph{Interval Computations} no.~2 (1993), 154--168.

\bibitem{hoffman1984} K.~L. Hoffman, G.~H. Parsons, L.~J. Allerdt, J.~M. Brooks and L.~E. Miles.
Elimination of ``hook-effect'' in two-site immunoradiometric assays by kinetic rate analysis.
\emph{Clinical Chemistry} 30 (1984), 1499--1501.
\href{https://doi.org/10.1093/clinchem/30.9.1499}{doi:10.1093/clinchem/30.9.1499}.

\bibitem{hoofnagle2009} A.~N. Hoofnagle and M.~H. Wener.
The fundamental flaws of immunoassays and potential solutions using tandem mass spectrometry.
\emph{Journal of Immunological Methods} 347 (2009), 3--11.
\href{https://doi.org/10.1016/j.jim.2009.06.003}{doi:10.1016/j.jim.2009.06.003}.

\bibitem{hunsaker2024} J.~J.~H. Hunsaker, S.~L. La'ulu, E. Zupan, D. Patel, V. Pandya and J.~W. Rudolf.
Detection of a high-dose hook effect and evaluation of dilutions of urine myoglobin specimens using a serum myoglobin assay.
\emph{Annals of Laboratory Medicine} 44 (2024), 367--370.
\href{https://doi.org/10.3343/alm.2023.0427}{doi:10.3343/alm.2023.0427}.

\bibitem{ichm10} International Council for Harmonisation.
\emph{ICH Harmonised Guideline M10: Bioanalytical Method Validation and Study Sample Analysis}.
Final version, adopted 24 May 2022.
\url{https://database.ich.org/sites/default/files/M10_Guideline_Step4_2022_0524.pdf}.

\bibitem{jacobs2015} J.~F.~M. Jacobs, R.~G. van~der Molen, X. Bossuyt and J. Damoiseaux.
Antigen excess in modern immunoassays: to anticipate on the unexpected.
\emph{Autoimmunity Reviews} 14 (2015), 160--167.
\href{https://doi.org/10.1016/j.autrev.2014.10.018}{doi:10.1016/j.autrev.2014.10.018}.

\bibitem{gum2008} Joint Committee for Guides in Metrology.
\emph{Evaluation of measurement data --- Guide to the expression of uncertainty in measurement}.
JCGM 100:2008, BIPM, 2008.
\href{https://doi.org/10.59161/JCGM100-2008E}{doi:10.59161/JCGM100-2008E}.

\bibitem{manski2003} C.~F. Manski.
\emph{Partial Identification of Probability Distributions}.
Springer Series in Statistics, Springer, New York, 2003.
\href{https://doi.org/10.1007/b97478}{doi:10.1007/b97478}.

\bibitem{mathlib2020} The mathlib Community.
The Lean mathematical library.
In \emph{Proceedings of the 9th ACM SIGPLAN International Conference on Certified Programs and Proofs (CPP 2020)}, ACM, 2020, pp.~367--381.
\href{https://doi.org/10.1145/3372885.3373824}{doi:10.1145/3372885.3373824}.

\bibitem{miles1968} L.~E.~M. Miles and C.~N. Hales.
Labelled antibodies and immunological assay systems.
\emph{Nature} 219 (1968), 186--189.
\href{https://doi.org/10.1038/219186a0}{doi:10.1038/219186a0}.

\bibitem{moore2009} R.~E. Moore, R.~B. Kearfott and M.~J. Cloud.
\emph{Introduction to Interval Analysis}.
SIAM, Philadelphia, 2009.
\href{https://doi.org/10.1137/1.9780898717716}{doi:10.1137/1.9780898717716}.

\bibitem{qian2003} S. Qian and H.~H. Bau.
A mathematical model of lateral flow bioreactions applied to sandwich assays.
\emph{Analytical Biochemistry} 322 (2003), 89--98.
\href{https://doi.org/10.1016/j.ab.2003.07.011}{doi:10.1016/j.ab.2003.07.011}.

\bibitem{rey2017} E.~G. Rey, D. O'Dell, S. Mehta and D. Erickson.
Mitigating the hook effect in lateral flow sandwich immunoassays using real-time reaction kinetics.
\emph{Analytical Chemistry} 89 (2017), 5095--5100.
\href{https://doi.org/10.1021/acs.analchem.7b00638}{doi:10.1021/acs.analchem.7b00638}.

\bibitem{rodbard1978a} D. Rodbard and Y. Feldman.
Kinetics of two-site immunoradiometric (`sandwich') assays---I. Mathematical models for simulation, optimization, and curve fitting.
\emph{Immunochemistry} 15 (1978), 71--76.
\href{https://doi.org/10.1016/0161-5890(78)90045-7}{doi:10.1016/0161-5890(78)90045-7}.

\bibitem{rodbard1978b} D. Rodbard, Y. Feldman, M.~L. Jaffe and L.~E.~M. Miles.
Kinetics of two-site immunoradiometric (`sandwich') assays---II. Studies on the nature of the `high-dose hook effect'.
\emph{Immunochemistry} 15 (1978), 77--82.
\href{https://doi.org/10.1016/0161-5890(78)90046-9}{doi:10.1016/0161-5890(78)90046-9}.

\bibitem{rodbardweiss1973} D. Rodbard and G.~H. Weiss.
Mathematical theory of immunoradiometric (labeled antibody) assays.
\emph{Analytical Biochemistry} 52 (1973), 10--44.
\href{https://doi.org/10.1016/0003-2697(73)90328-X}{doi:10.1016/0003-2697(73)90328-X}.

\bibitem{ross2020} G.~M.~S. Ross, D. Filippini, M.~W.~F. Nielen and G.~IJ. Salentijn.
Unraveling the hook effect: a comprehensive study of high antigen concentration effects in sandwich lateral flow immunoassays.
\emph{Analytical Chemistry} 92 (2020), 15587--15595.
\href{https://doi.org/10.1021/acs.analchem.0c03740}{doi:10.1021/acs.analchem.0c03740}.

\bibitem{roy2017} R.~D. Roy, C. Rosenmund and M.~I. Stefan.
Cooperative binding mitigates the high-dose hook effect.
\emph{BMC Systems Biology} 11 (2017), 74.
\href{https://doi.org/10.1186/s12918-017-0447-8}{doi:10.1186/s12918-017-0447-8}.

\bibitem{ryall1982} R.~G. Ryall, C.~J. Story and D.~R. Turner.
Reappraisal of the causes of the ``hook effect'' in two-site immunoradiometric assays.
\emph{Analytical Biochemistry} 127 (1982), 308--315.
\href{https://doi.org/10.1016/0003-2697(82)90178-6}{doi:10.1016/0003-2697(82)90178-6}.

\bibitem{sadler1988} W.~A. Sadler, M.~H. Smith and H.~M. Legge.
A method for direct estimation of imprecision profiles, with reference to immunoassay data.
\emph{Clinical Chemistry} 34 (1988), 1058--1061.
\href{https://doi.org/10.1093/clinchem/34.6.1058}{doi:10.1093/clinchem/34.6.1058}.

\bibitem{sathishkumar2020} N. Sathishkumar and B.~J. Toley.
Development of an experimental method to overcome the hook effect in sandwich-type lateral flow immunoassays guided by computational modelling.
\emph{Sensors and Actuators B: Chemical} 324 (2020), 128756.
\href{https://doi.org/10.1016/j.snb.2020.128756}{doi:10.1016/j.snb.2020.128756}.

\bibitem{selby1999} C. Selby.
Interference in immunoassay.
\emph{Annals of Clinical Biochemistry} 36 (1999), 704--721.
\href{https://doi.org/10.1177/000456329903600603}{doi:10.1177/000456329903600603}.

\bibitem{sturgeon2011} C.~M. Sturgeon and A. Viljoen.
Analytical error and interference in immunoassay: minimizing risk.
\emph{Annals of Clinical Biochemistry} 48 (2011), 418--432.
\href{https://doi.org/10.1258/acb.2011.011073}{doi:10.1258/acb.2011.011073}.

\bibitem{tate2004} J. Tate and G. Ward.
Interferences in immunoassay.
\emph{The Clinical Biochemist Reviews} 25 (2004), 105--120.

\bibitem{wild2013} D. Wild (ed.).
\emph{The Immunoassay Handbook: Theory and Applications of Ligand Binding, ELISA and Related Techniques}, 4th ed.
Elsevier, Oxford, 2013.

\bibitem{wu2018} S.~J. Wu and J.~A. Hayden.
Upfront dilution of ferritin samples to reduce hook effect, improve turnaround time and reduce costs.
\emph{Biochemia Medica} 28 (2018), 010903.
\href{https://doi.org/10.11613/BM.2018.010903}{doi:10.11613/BM.2018.010903}.
